\section{Introduction}\label{sec:introduction}

Gr\"atzer and Schmidt proved that every algebraic lattice is the full
congruence lattice of an algebra \citep[Theorem~I]{GratzerSchmidt1963}.
The carrier in that theorem need not be finite.  The finite lattice
representation problem asks whether every finite lattice is
the congruence lattice of a finite algebra.  The operations of the algebra are
unrestricted: their number, arities, and tables may depend on the lattice.
Thus a negative answer must exclude all finite algebras, rather than only
groups, rings, or algebras in a fixed variety.

There is also a decision question: given the finite order table of a
lattice, can one determine whether such a finite representation exists?
McNulty records this precise question separately from the universal
representation problem \citep[p.~30]{McNulty2015}.  We prove that no such
decision procedure exists, and hence that not every finite lattice has a
finite representation.  We also construct a nonrepresentable family
directly.  Both conclusions must be distinguished from the effective
search for affirmative witnesses, which is possible.

The relation with finite groups is particularly close.
P\'alfy and Pudl\'ak proved that every finite lattice is a congruence lattice
of a finite algebra if and only if every finite lattice is an interval in a
subgroup lattice of a finite group \citep[Theorem~2]{PalfyPudlak1980}.
This is an equivalence between two universally quantified statements.
It does not assert that these two kinds of representation are equivalent
for each individual lattice.  Our reductions below are proved for the
particular lattice families to which they are applied.

Finite partition embeddings do not settle the problem: every finite
lattice embeds in the equivalence lattice of a finite set
\citep[Theorem~7.10]{PudlakTuma1980}, but the image need not be the full
congruence lattice of an algebra on that set.  Positive representation
results include all finite lattices of width at most two
\citep[Theorem~4.2]{Snow2000}.  Small examples likewise give little
indication of the obstruction.
DeMeo's analysis leaves only one possible exception among lattices with at
most seven elements \citep[Theorem~6.1.1 and Figure~6.2]{DeMeo2012}.
We represent that exception explicitly in \Cref{sec:small}.
The obstruction in this paper instead uses long systems of affine
subgeometries and the restrictions that their subgroup labels impose
simultaneously in several ambient simple groups.

\subsection{Statements}

A finite algebra means a nonempty finite set \(A\), together with a finite
list of total operations of finite nonnegative arity.  The list may be
empty.  An operation of arity zero specifies a constant and imposes no
additional condition on an equivalence relation.
Write \(\Con(A)\) for the lattice of all congruences of the algebra.
For a nonempty finite lattice \(L\), let
\[
 \mathcal R(L)\quad\Longleftrightarrow\quad
 L\cong\Con(A)\text{ for some finite algebra }A .
\]
Every representation in this paper is onto the entire congruence lattice.
An embedding into a proper sublattice is not a representation in this sense.

\Needspace{8\baselineskip}
\begin{theorem}\label{m:main}
The following statements hold.
\begin{enumerate}[label=\textup{(\roman*)}]
\item A finite nonempty lattice is representable if and only if it satisfies
the finite colored-graph criterion in \Cref{f:criterion}.
For each candidate graph size the criterion is decidable by a finite test.
\item There is no deterministic Turing machine that, given the order table
of a finite nonempty lattice, always halts and decides \(\mathcal R(L)\).
\item There is a nonrepresentable finite lattice \(L_0\) of minimum
cardinality, specified by the finite formulas in \Cref{sec:minimum}.
Every finite lattice of smaller cardinality is representable, and
\(\lvert L_0\rvert>7\).
\item Given the order table of a finite nonempty lattice \(L\), it is
undecidable whether \(L\cong[H,G]\) for some finite group \(G\) and
subgroup \(H\le G\), where \([H,G]\) is the full subgroup interval.
\end{enumerate}
\end{theorem}

Part~\textup{(ii)} resolves the finite lattice representation problem
negatively, since otherwise the constant affirmative algorithm would be a
decider.  We also construct a nonrepresentable family directly, before the
undecidability argument.  Part~\textup{(iii)} specifies an actual finite
order table by an unambiguous finite description.  Its constants include
one ordinary Busy Beaver integer.  The description does not evaluate
that integer, the minimum cardinality, or the entries of the order table.

The distinction between the three congruence-lattice assertions is useful.
The criterion
in Part~\textup{(i)} gives a search that terminates on representable
lattices; it does not provide a computable stopping rule for negative
instances.  The proof of Part~\textup{(ii)} establishes precisely that
there is no uniform such rule.  In Part~\textup{(iii)}, a single finite
search range is specified mathematically, using the Busy Beaver
constant; no uniform algorithm for evaluating these constants is asserted.

\subsection{Main ideas and proof organization}

We first make the affirmative search explicit.  Edge colors of a finite
complete graph encode principal congruences, and label-decreasing
selfmaps replace the choice of an algebraic signature.  A path condition
ensures that the recovered relations exhaust the congruence lattice.
For fixed graph size this is a finite test.  Enumerating graphs therefore
finds every affirmative witness, and a decider exists if and only if a
computable bound on witness sizes exists (\Cref{sec:graph}).  This is the
effective bridge used in the undecidability proof.

The graph method itself is established background.  Finite unary closure
already gives positive semidecidability
\citep[Section~3.2, Theorem~3.2.1]{DeMeo2012}.  Pudl\'ak's
valuation-colored construction uses graph paths and stable maps, which
preserve an edge color or collapse the edge; its general representation
can have an infinite carrier \citep[pp.~270--272]{Pudlak1976}.
Kenney's graphical-algebra framework already uses label-decreasing maps,
principal-congruence complete graphs, edge-image paths, and reconstruction
by all unary graph endomorphisms
\citep{Kenney2010}; see also \citet[pp.~51--58]{Kenney2009Slides}.
We present the finite criterion directly in terms of the lattice and one
finite graph, with the proof and bounded arithmetic needed for the final
minimum specification.

The first structural construction produces a negative family, without
using a hypothetical decider.  We assemble a simple self-dual lattice
from private truncated affine charts.  On a least carrier, every unary
operation is a permutation or a constant, but the permutation action
need not be transitive.  A full-diagonal construction in a semidirect
power represents the dual lattice as an ordinary subgroup interval;
self-duality gives the required orientation.  Normal-subgroup tests then
leave either an almost simple interval or extension-kernel data.  A
four-dimensional affine grid excludes the latter, yielding an ordinary
almost simple interval for the whole lattice or an initial point link
(\Cref{sec:affine}).

The ordinary-interval step develops methods used by P\'alfy and
Pudl\'ak, including restriction along idempotents on a least carrier and,
in the later subdirect descent, the correspondence between transitive
actions and stabilizer intervals
\citep[Lemmas~1 and~3]{PalfyPudlak1980}.  The subsequent extension
analysis is related to the twisted-wreath and signalizer approaches of
Baddeley and Lucchini, Baddeley, B\"orner, and Aschbacher
\citep{BaddeleyLucchini1997,Baddeley1998,Borner1999,Aschbacher2008};
P\'alfy compares their hypotheses and interval reductions
\citep[Section~4]{PalfyTwisted}.  These results concern different classes
of lattices and do not supply a pointwise reduction for the present
family.  Here the normal tests, compatibility of actual simple factors,
and extension-kernel grid provide that reduction directly.

Recognizing abstract simple-group types at consecutive labels would not
suffice for the obstruction.  The affine patterns force actual nested
simple subgroups, while exclusions are imposed relative to every earlier
ancestor at once (\Cref{sec:patterns,sec:classical,sec:alternating}).
Bounded Lie rank is excluded by algebraic subgroup tests and the
sufficiently-general subgroup method of Larsen and Pink
\citep[Theorem~0.5]{LarsenPink2011}; \Cref{sec:quantitative} supplies an
effective version for the later size ceiling.  In high rank, one must
further exclude field roots and tensor decompositions in every ancestor.
These protections make the nested actions compatible with rational
composition.  Highest-weight theory then produces universal cuts in the
full weight multiset, and multiplicities rule out a sufficiently long
protected path (\Cref{sec:bounded,sec:weights}).  Choosing more affine
generations than this bound gives the nonrepresentable family.

Undecidability requires a different uniformity statement.  A weighted
template is a subgroup interval built from finitely many permutation
sets, called units, whose cardinalities encode integer values.  A
positive solution of an arithmetic circuit equation supplies such a
template, already finitely representable by the coset construction.
The number of units depends only on the circuit's syntax, not its
solution.  We prove that any finite representation yields a new reading
of all the encoded values in one common partition into cells.  Equal
totals encode addition, and square cardinalities encode multiplication;
the new values need not equal the defining ones.

\Needspace{3\baselineskip}
Two independent bounds make this useful.  First, cell sizes are bounded
computably in the size of a representation.  Second, the number of
elements of a template is bounded computably in its number of units,
independently of their weights.  A hypothetical decider would thus bound
the sizes of affirmative witnesses for every template arising from a
fixed circuit.  Reading such a bounded witness would replace any
positive solution by a computably bounded one.  Searching the resulting
finite box contradicts Diophantine undecidability.  No algorithm is
required to construct the unknown-weight template from the circuit
(\Cref{sec:counting,sec:diophantine}).  The same replacement argument,
using bounds on top-group orders instead of algebra carriers, proves
undecidability of full subgroup-interval recognition.

The common reading is the structural burden of this second branch.
Reverse alternating tests reuse the pruning argument through
elementary-linear charts.  A forward marker theorem instead uses marker
omissions and their three-coatom constraints to reach the protected-path
bound (\Cref{sec:reverse,sec:forward,sec:forward-completion}).  Together
with the weighted normal tests, they force the
remaining units into an actual alternating power.  Normal and parity
tests then recover one cell alignment in which every total and square
comparison holds, even after the nontrivial lattice images needed by the
least-carrier reduction (\Cref{sec:weighted}).  The separate counting
argument concerns the literal templates, not those readings.  It uses
nonabelian coordinate compression, effective integer profiles,
bounded-tuple encodings of diagonal matchings, and Mar\'oti's
primitive-group order theorem \citep[Theorem~1.1]{Maroti2002}.

Finally, \Cref{sec:quantitative} gives an ordinary computable size ceiling
for one member of the nonrepresentable affine family.  This step is
independent of both the hypothetical decider and the later
Busy Beaver constant.  The latter constant bounds the witness sizes
returned by the prescribed affirmative graph searches up to the ceiling.
Explicit bounded arithmetic then
selects the first negative order table, ordered first by cardinality
(\Cref{sec:minimum}).

\begin{figure}[tbp]
\centering
\begin{tikzpicture}[
  box/.style={draw,rounded corners=2pt,align=center,text width=4.7cm,
    inner sep=6pt,font=\small},
  arr/.style={-{Stealth[length=2mm]},thick},
  node distance=10mm and 13mm]
\node[box,text width=8.2cm] (engine)
  {Simultaneous pruning in actual simple groups\\
   and the protected-path bound};
\node[box] (affine) at (-3.2,-1.7)
  {Private affine charts:\\a nonrepresentable family};
\node[box] (marker) at (3.2,-1.7)
  {Reverse tests and\\the forward marker theorem};
\node[box,below=of affine] (ceiling)
  {Effective parameter control:\\a computable size ceiling\\for one negative lattice};
\node[box,below=of marker] (weighted)
  {Weighted representations:\\one cell alignment for\\all arithmetic readings};
\node[box,below=of ceiling] (graph)
  {Finite graph witnesses:\\affirmative search and\\conditional size bounds};
\node[box,below=of weighted] (count)
  {Independent template count:\\lattice size bounded by units,\\not numerical weights};
\node[box,below=of graph] (minimum)
  {An ordinary Busy Beaver cutoff\\and finite selection of a\\minimum negative order table};
\node[box,below=of count] (undec)
  {Bounded replacement of\\Diophantine solutions:\\no representation decider};
\draw[arr] (engine.south west) -- (affine.north);
\draw[arr] (engine.south east) -- (marker.north);
\draw[arr] (affine) -- (ceiling);
\draw[arr] (marker) -- (weighted);
\draw[arr] (ceiling.west) -- ++(-0.35,0) |- (minimum.west);
\draw[arr] (graph) -- (minimum);
\draw[arr] (graph.south east) -- (undec.north west);
\draw[arr] (count) -- (undec);
\draw[arr] (weighted.east) -- ++(0.35,0) |- (undec.east);
\end{tikzpicture}
\caption{The two uses of the structural argument.  The affine branch
gives a negative lattice and a computable size ceiling independently of
any hypothetical decider.  The weighted branch combines simultaneous
arithmetic readings with a separate, weight-independent count; a
decider would then bound positive Diophantine solutions.  Finite graph
witnesses also supply the bounded searches used, with the size ceiling,
in the minimum-table specification.}
\label{fig:proof-overview}
\end{figure}


\subsection{Conventions and dependencies}

All finite lattices and algebra carriers are nonempty.  The one-element
lattice is included.  A subgroup interval \([H,G]\) always means the
whole interval.  Invariant-subgroup intervals in a normal direct power
will be identified explicitly and are not silently substituted for
ordinary intervals.  The notation \(\core_G(H)\) denotes the largest
normal subgroup of \(G\) contained in \(H\).
Quotienting by a normal subgroup already contained in the base preserves
the full interval above the base.

Labels such as \(K_X\) always denote actual subgroups in the realization
or quotient currently under discussion.  In a reverse affine chart,
a smaller flat has a larger subgroup label.
A normal subgroup is \emph{retained} when it lies in the current base;
it is \emph{missing} when the relevant normal test forces the base
to supplement it.  An element \(y\) of a lattice is left-modular if
\[
 x\vee(y\wedge z)=(x\vee y)\wedge z
 \qquad\text{whenever }x\leq z.
\]
The field of an affine chart and the field of a group's natural module
need not be the same.  Group bases, field letters and local ranks are
reintroduced when their ambient objects change.

We use the classification and standard structure of finite simple
groups, algebraic groups and their representations.  The exact
conventions and applicable ranges are recorded in
\Cref{sec:affine} and at their later uses.
Established results are cited; the structural tests, compatibility
arguments, weighted readings and effectivity deductions specific to
the constructions are proved below.
\FloatBarrier
